% D8.1 The four conditions of E4 and the comparison each step makes.
% Shared figure colours (colour only in figures).
\definecolor{robotfill}{HTML}{EAF2FC}\definecolor{robotline}{HTML}{2A78D6}
\definecolor{serverfill}{HTML}{FDEEE7}\definecolor{serverline}{HTML}{EB6834}
\definecolor{netfill}{HTML}{F0EFEC}\definecolor{greenfill}{HTML}{E9F5EC}\definecolor{greenline}{HTML}{2F9E58}
\definecolor{purplefill}{HTML}{F1ECF9}\definecolor{purpleline}{HTML}{7A55C7}
\definecolor{ink}{HTML}{0B0B0B}\definecolor{ink2}{HTML}{52514E}\definecolor{faultc}{HTML}{C9C6BF}

\begin{tikzpicture}[font=\footnotesize, node distance=5mm and 19mm,
  st/.style={draw=#1, thick, rounded corners=3pt, fill=#1!8, align=center, text width=31mm, minimum height=17mm},
  cmp/.style={-{Latex[length=2.2mm]}, very thick, ink2}, lab/.style={font=\scriptsize, text=ink2, align=center}]
\node[st=robotline] (c1) {{\bfseries 1. Reference}\\ patient LeRobot,\\ normal network};
\node[st=serverline, right=of c1] (c2) {{\bfseries 2. Faults}\\ patient LeRobot,\\ the fault mix};
\node[st=greenline, right=of c2] (c3) {{\bfseries 3. \drtc}\\ \drtc's scheduling,\\ the fault mix};
\node[st=purpleline, right=of c3] (c4) {{\bfseries 4. \drtc{} + stack}\\ adaptive, ensembling,\\ Butterworth, ramp; the mix};
\draw[cmp] (c1) -- node[lab, above, text width=16mm]{what the faults cost} (c2);
\draw[cmp] (c2) -- node[lab, above, text width=16mm]{what \drtc{} wins back} (c3);
\draw[cmp] (c3) -- node[lab, above, text width=16mm]{what smoothing adds} (c4);
\node[lab, below=4mm of c2.south east] {primary comparisons: 2 vs 1, 3 vs 2, 4 vs 3, each with M1 (task success);
  4 policies $\times$ 10 trials per condition on cube in box v1, interleaved in seeded rounds};
\end{tikzpicture}
